% ED26-BM-RAT2021-5-5: exact 2021 numbered chapter, pp. 361--366.
\chapter{Argument Strength, Dialogue, and Probabilistic Assessment}
\label{sec:p18-rationality-5-5}

Being persuaded, having a good reason, and belonging to an accepted
argument extension are different outcomes. Hahn and Collins's 2021
chapter asks both what convinces people and what ought to convince them.
Its account links argument content, dialogue, normative justification
and empirical evaluation; it does not reduce argument quality to the
attack relation of an abstract framework
\cite[pp.~361--364]{HahnCollins2021ArgumentationTheory}.
This treatment retains that perspective and supplies explicit finite
calculations for its probabilistic illustrations. The numerical models
below are stipulated explanatory examples, not medical evidence or
participant data reported in the source.

\section{Arguments as objects and as processes}

\begin{hbdefinition}{def:p18-r55-objects-processes}{Two objects of assessment}
An argument considered as an \emph{object} consists of a conclusion
and premises offered in its support. The chapter notes that some uses
of this expression refer just to the offered reasons; the present
convention includes their claimed conclusion. A deductive argument
purports to exclude true premises with a false conclusion. An ampliative
argument can have true premises and a false conclusion: statistical,
abductive and defeasible arguments fall within the chapter's deliberately
broad use of \emph{inductive}. Their differences are not thereby erased.
Defeasible endorsement is provisional and may be overturned by further
information; that qualification is part of the offered conclusion.

An argument considered as a \emph{process} is an exchange in a social
setting: participants advance, challenge, defend or withdraw claims,
or propose actions. The content can be reconstructed as an object,
but the exchange additionally has commitments, purposes, utterance
contexts and obligations to respond. A report from a speaker is not
automatically knowledge that the reported proposition is true.
\cite[p.~361 and note~1, p.~365]{HahnCollins2021ArgumentationTheory}
\index{argumentation!object and process}\index{argument strength}
\end{hbdefinition}

Concentrating on objects alone can miss why a premise would be uttered,
what information its speaker possesses, and what a listener is entitled
to infer from the utterance. The source uses deviations from modus
tollens as a reason to investigate conversational context, rather than
declaring either every deviation irrational or every contextual reading
deductively valid. Conversely, concentrating only on dialogue rules
can leave the evidential quality of the exchanged reasons unexplained.
Neither vantage point replaces the other.

\begin{hbdefinition}{def:p18-r55-norms}{Procedural, epistemic and conventional standards}
\emph{Procedural norms} govern participants' conduct: advancing a
standpoint, securing commitments, responding to challenges, or meeting
a burden of proof. Their purposes can include resolving disagreement,
fair treatment and orderly participation. \emph{Epistemic norms}
assess the content and inferential relationship of reasons to conclusions,
with truth or justified belief as aims. Logical consistency and
probabilistic coherence are examples, not identical tests.

\emph{Conventional validity}, as used in the chapter, is an empirical
question about whether ordinary discussants' judgments or practices
accord with proposed norms. Agreement with a norm is different from
a justification of why the norm should be followed. Judgments matching
its predictions are also different from evidence that participants
actually used that norm in producing those judgments.
\cite[pp.~362--363]{HahnCollins2021ArgumentationTheory}
\end{hbdefinition}

\section{The normative argument and its objections}
\label{sec:p18-r55-normative-argument}

The source discusses pragma-dialectics as a prominent procedural account:
different stages of a discussion have different rules, including an
obligation to offer support for an asserted standpoint. Its epistemic
comparison starts from classical validity and then considers probability
as a way to assess ampliative support. A conclusion can receive stronger
support without becoming a necessary consequence or being certain.

The source's circular example asserts a deity's existence using that
same existence as its reason. Writing \(H\) for that proposition here,
its deductive form \(H\vdash H\)
is valid; that fact supplies no independent reason for an interlocutor
who doubts \(H\). A procedural explanation might say that the speaker
has failed to discharge a burden. Hahn and Collins object that explaining
the evidential weakness solely by that failure can presuppose the very
weakness being explained: one still needs a content-sensitive account
of what would supply support. This is their argument for considering
epistemic and procedural norms together, not a theorem that every
dialogue obligation is epistemically redundant
\cite[pp.~361--362]{HahnCollins2021ArgumentationTheory}.

Their Bayesian perspective assesses how offered information should alter
degrees of belief in a conclusion. An increase is a positive evidential
contribution; its magnitude, final confidence and adequacy for a decision
are further questions. The chapter also allows utilities when an argument
supports an action. Its discussion recognizes objections to probabilistic
models, competing predictions about ad hominem arguments, ethical
procedural goals, and virtue-based accounts that cross the content/process
distinction. It proposes complementarity, not the replacement of every
procedural or ethical standard by a posterior probability.

The authors appeal to accuracy-based mathematical foundations of Bayesian
inference, while treating detailed foundations of procedural norms as an
unfinished issue in their 2021 discussion. Their chapter does not state
one universal theorem covering every loss function, decision aim and
information acquisition process. The finite quadratic-score result below
is an explicit conditional justification, not an attribution of such a
universal theorem to the cited literature. Whether consensus rules promote
truth, fairness, participation or several of these goals requires the
appropriate separate argument and evidence.

\section{The toxic-effect illustration: evidence and its absence}

The source contrasts support for a toxicity claim from an observed effect
with support for nontoxicity from no observed effect. The second is an
argument from ignorance, but its quality depends on how likely a toxic
effect would have been detected. Calling it an instance of that scheme
does not settle whether it supplies evidence
\cite[p.~363]{HahnCollins2021ArgumentationTheory}.

\begin{hbdefinition}{def:p18-r55-detection-model}{A stated detection experiment}
In a finite probability model let \(H\) denote toxicity and \(E\)
an observed toxic effect. Let \(H^c,E^c\) be their complements and set
\[
 \pi=\mathbb P(H),\qquad
 s=\mathbb P(E\mid H),\qquad
 c=\mathbb P(E^c\mid H^c),\qquad 0<\pi<1.
\]
Here \(s\) is sensitivity and \(c\) specificity; neither is a
posterior probability. Conditioning presumes the stated observation
event has positive probability and that the stipulated likelihoods
describe the relevant experiment. A failure to perform a test is not
the same evidence as a specified test with no detected effect.
\end{hbdefinition}

\begin{hbtheorem}{thm:p18-r55-detection}{proved=annex:p18.r55.detection}{Positive and negative evidential force}
When their conditioning events have positive probability,
\[
\begin{split}
 p_+&=\mathbb P(H\mid E)
      =\frac{\pi s}{\pi s+(1-\pi)(1-c)},\\
 q_-&=\mathbb P(H^c\mid E^c)
      =\frac{(1-\pi)c}{\pi(1-s)+(1-\pi)c}.
\end{split}
\]
The observed effect raises \(H\)'s probability above \(\pi\),
and the observed absence raises \(H^c\)'s probability above \(1-\pi\),
exactly when \(s+c>1\). Equality gives no change; \(s+c<1\)
reverses these directions. These are statements about the two different
conclusions under the given joint model, not a diagnostic recommendation.
\end{hbtheorem}

\begin{hbexample}{def:p18-r55-detection-values}{Strong absence and high confidence without support}
For the stipulated values \(\pi=1/10\), \(s=4/5\), \(c=9/10\),
the detected-effect posterior is \(p_+=8/17\), while the no-effect
posterior for nontoxicity is \(q_-=81/83\). Neither is deduced solely
from the argument's verbal form. With the same prior but
\(s=1/50\), \(c=49/50\), the observation is independent of
toxicity: \(p_+=1/10\) and \(q_-=9/10\). The high latter
probability was already the prior; the absence has supplied no increase.
A statement of posterior confidence must therefore not be substituted
for a statement of evidential improvement.
\end{hbexample}

The probability foundation is the actual conditioning convention and
finite Bayes result at
\cref{def:p15-finite-conditioning,thm:p15-finite-bayes}.
Calibration of \(s,c,\pi\), selection of tested cases, dependence
between observations, and changes in the experiment require their own
evidence. An observed proportion or an estimated likelihood does not
become an exact probability by being put into the displayed formula.

\section{Accuracy, duplicated reasons, action and testimony}

\begin{hbtheorem}{thm:p18-r55-norm-boundaries}{proved=annex:p18.r55.norm-boundaries}{Four distinct conditional standards}
The following finite facts delimit the preceding normative discussion.
\begin{enumerate}
\item If an evidence event \(E\) entails a conclusion event \(H\),
      meaning \(E\subseteq H\), and \(\mathbb P(E)>0\), then
      \(\mathbb P(H\mid E)=1\). A posterior of one is not by
      itself proof of logical entailment: probability-zero countercases
      can remain in the model.
\item For a Bernoulli truth indicator \(h\), a probability report
      \(r\in[0,1]\), and quadratic loss \((r-h)^2\), expected
      loss under conditional probability \(p\) is
      \((r-p)^2+p(1-p)\), uniquely minimized at \(r=p\).
      This conclusion assumes that loss and that conditional model.
\item If two report events are identical, \(E_2=E_1\), and
      \(\mathbb P(E_1)>0\), then
      \(\mathbb P(H\mid E_1\cap E_2)=\mathbb P(H\mid E_1)\).
      Counting the copied report again supplies no new evidence.
\item In a stipulated action comparison with utility \(10\) for a
      favorable outcome, \(-100\) for an unfavorable outcome of
      probability \(q\), and utility zero for the alternative,
      expected-utility choice favors the first action exactly when
      \(q<1/11\), with a tie at equality.
\end{enumerate}
The last criterion is not a posterior-confidence threshold independent
of the utilities. The score-minimization result does not prove that
every participant uses quadratic loss or that every ethical dialogue
goal is reducible to that loss.
\end{hbtheorem}

The duplicated-report calculation is one precise limitation of treating
repetition as additional support. It is not a definition covering every
kind of circular argument. Likewise a likelihood product for multiple
reports requires an appropriate conditional-independence model, not
merely different speakers or different words.

The source's slippery-slope illustration connects allowing same-sex
marriage with an asserted later demand for pet marriage. It is an
illustrative argument, not a verified causal progression. The proposed
probabilistic analysis distinguishes the probability of an outcome
given an action from its utility. Its reported experimental connection
to categorization concerns how such arguments are assessed, not whether
the predicted social outcome occurs. The explicit utility calculation
above shows why changing a consequence's cost can change an action
assessment even when its probability is held fixed
\cite[p.~363]{HahnCollins2021ArgumentationTheory}.

For testimony, the relevant event is that a source makes a report,
not simply the report's propositional content being true. Expertise,
honesty, incentives, access to information and dependence between
sources can alter the likelihood of that report under competing
hypotheses. A source criticism can therefore be evidentially relevant;
an abusive remark unrelated to those likelihoods supplies no such
probabilistic account. The chapter reports Bayesian models of expertise
and ad hominem assessment, without claiming that any personal attack
is justified or that a person's conclusion is false merely because
the person is disliked. The schemes and critical-question treatment
at \cref{sec:p18-argument-schemes-critical-questions} supplies the
distinct reconstruction task: identify what is being claimed about
the source and whether evidence actually supports that criticism.

\section{What the empirical evidence does and does not establish}

Hahn and Collins organize the evidence into naturalistic analyses,
controlled argument judgments, individual differences and intervention
research. Their examples include procedural analyses of extended
discussions, probabilistic reconstruction of the Sacco and Vanzetti
evidence, and taxonomies containing more than fifty schemes. These
are identified research approaches, not three interchangeable estimates
of argument quality. They also note possible redundancy in schemes
and incomplete coverage, especially of causal argument. A taxonomy's
length neither proves its completeness nor supplies probabilities.

Their experiment survey describes judgments of reasonableness,
persuasiveness or convincingness, often on rating scales. Reported
procedural comparisons distinguish legitimate personal criticism,
disguised fallacious criticism and open abuse. Scheme studies examine
authority, analogy and examples. Bayesian studies vary sensitivity,
specificity, consequence probabilities/utilities or source properties,
often in dialogues. In the chapter these are qualitative reports:
the individual study samples, effect estimates and uncertainty are not
supplied. They are not population frequencies, individual-level fits,
or evidence that one framework uniquely explains all the judgments
\cite[p.~363]{HahnCollins2021ArgumentationTheory}.

One of the chapter's cited primary studies permits a more definite
bounded record. Schellens and colleagues interviewed 108 untrained
participants: 24 Dutch university students for authority, and 84 library
visitors divided among four other schemes. Participants generated
strong/weak arguments, ranked supplied arguments and explained their
reasons; claim acceptance was not measured. A second rater coded
115 protocol segments, with reported \(\kappa\) from \(0.69\)
to \(0.80\). In the authority group, expertise appeared in 24/24
participants' protocols and trustworthiness in 22/24. These are task-bound
coding counts, not population proportions or calibrated probabilities.
Recruitment was by the indicated university and library settings, not a
population probability sample. The authors acknowledge material/task effects and possible interpretive
coding bias; they decline a conclusion about relative frequencies of
general versus scheme-specific criteria. No confidence interval for
population prevalence is supplied. The result supports use of such
criteria in these elicited tasks, not proof of their universal normative
correctness or everyday prevalence
\cite[pp.~686--689, 697--700]{SchellensEtAl2017SchemeCriteria}.

The 2021 chapter's group-level conclusion is accordingly limited:
there is evidence of sensitivity to procedural, scheme-specific and
probabilistic considerations. Such sensitivity is compatible with
rational assessment in the tested settings. It does not establish
uniform competence, uniquely identify the underlying computation,
or remove uncertainty about which norm is justified for which purpose.
An empirical association with an endorsed norm and a conceptual
justification of that norm are different claims.

\section{Comparison with abstract and structured argumentation}

In the existing abstract treatment, a framework's supplied objects are
argument identifiers and attacks. Acceptability relative to a set means
defense against attackers. Grounded, complete, preferred and stable
extensions impose different conditions on that graph. Those definitions
do not include a sampling experiment, prior, likelihood or utility
(see \cref{def:p18-finite-framework,def:p18-defense-operator,def:p18-extension-semantics}).

\begin{hbtheorem}{thm:p18-r55-graph-boundary}{proved=annex:p18.r55.graph-boundary}{Extension acceptance does not identify probabilistic force}
For the attack-free one-argument framework \(\mathcal F=(\{a\},\varnothing)\),
the unique grounded, complete, preferred and stable extension is \(\{a\}\).
Associate \(a\) with the inference from a reported effect \(E\)
to \(H\). Two coherent finite probability models can have this same
framework and the same prior \(\mathbb P(H)=1/2\), but respectively
\[
 \mathbb P(H\mid E)=3/4\qquad\text{and}\qquad
 \mathbb P(H\mid E)=1/4.
\]
In the first the observation supports \(H\); in the second it
disconfirms \(H\). No probability-strength or confirmation verdict
is determined by that common extension verdict alone.
\end{hbtheorem}

This is an editorial separation result, not a claim that Dung intended
his semantics to estimate those probabilities or that Bayesian assessment
replaces graph-based analysis. A probabilistic augmentation of an attack
framework would need to specify its random objects, dependence model,
interpretation of attacks and aggregation/update rule; none follows
merely from labeling an argument accepted.

The structured treatment additionally specifies premises, strict and
defeasible rules, subarguments, contraries and preferences. Its undermining,
rebutting and undercutting relations identify where challenges reach an
argument; some attacks defeat regardless of preference, others use a
comparison at the targeted subargument
(see \cref{def:p18-aspic-system,def:p18-aspic-attacks,def:p18-aspic-defeat}).
An order of preference there is not automatically a conditional probability,
and its translation to extension semantics does not supply a likelihood
model (see \cref{def:p18-aspic-semantics}). The approaches can contribute
different layers of one analysis only after their linking assumptions
are made explicit: reconstruction, content-sensitive uncertainty,
challenge/defeat, and the commitments of the dialogue.

\section{Research directions and a usable assessment record}

The source's 2021 research agenda asks for individual-difference studies,
reliable measures of skill, interventions and tests of dispute resolution.
It warns that measuring skill by counts of rebuttals or counterarguments
can privilege procedural performance while leaving content quality poorly
assessed. The review relates education and epistemological beliefs to
success or participation, including its contrast between
\emph{multiplists}, who treat facts as freely chosen opinions, and
\emph{evaluativists}, who treat judgments as answerable to evidence.
Those associations do not by themselves establish a causal effect or
a complete theory of individual rationality.

The intervention reports span a three-year program with children and
an hourlong intervention with undergraduates. Reported transfer examples
include mathematical-argumentation training followed by mathematics and
reading standardized-test performance. The
chapter does not supply the individual samples, effect sizes, uncertainty
or a common design for those reports. It treats evidence as suggestive,
and describes a lack of equivalent scheme-based or Bayesian-norm
interventions as an open issue in its review, not as their impossibility.
Its proposed contrastive studies would need a defensible measure of skill
and a way to distinguish training-specific performance from transfer.

Finally it proposes adapting group judgment methods to ask whether
particular norms improve the probability or efficiency of consensus,
and whether that consensus tracks truth. Reaching agreement and reaching
an accurate answer must be measured separately. These are proposed
investigations, not results established by the chapter. Constrained
probabilistic models offer discriminating predictions, but only when
their parameters and response link are fixed independently enough to
permit failure; unconstrained post-hoc adjustment can conceal competing
explanations. The chapter also treats pragmatic implications as an area
where particular argument models may sharpen general language theories
\cite[p.~364]{HahnCollins2021ArgumentationTheory}.

\begin{hbmethod}{def:p18-r55-assessment-record}{Keep the queries separate}
For an actual argument, record its conclusion, offered reasons and
dialogue purpose. Separate an asserted premise from independent evidence
for it. Identify the claimed force: deductive consequence, positive
probabilistic support, adequate posterior confidence, or an action choice.
For a probability claim give the prior, observation model and dependence
assumptions; for an action claim also give the utilities and alternatives.
Ask the source-specific critical questions about these assumptions.

If an attack framework is also used, give its supplied argument/attack
construction and selected extension semantics, reporting that output
separately. Check whether new information is genuinely new, whether an
absence could have been detected, and whether a personal criticism changes
the evidential model. For an empirical claim retain the actual task,
sample, response measure, uncertainty and competing explanation.
End with a qualified conclusion and the remaining unanswered question;
neither a fallacy label, a favorable rating nor an extension membership
can substitute for those missing inputs.
\end{hbmethod}
